% Options for packages loaded elsewhere
\PassOptionsToPackage{unicode}{hyperref}
\PassOptionsToPackage{hyphens}{url}
\PassOptionsToPackage{dvipsnames,svgnames,x11names}{xcolor}
\documentclass[
  10pt,
]{article}
\usepackage{xcolor}
\usepackage[margin=1cm,top=1cm,bottom=1cm,left=1cm,right=1cm,includeheadfoot]{geometry}
\usepackage{amsmath,amssymb}
\setcounter{secnumdepth}{3}
\usepackage{iftex}
\ifPDFTeX
  \usepackage[T1]{fontenc}
  \usepackage[utf8]{inputenc}
  \usepackage{textcomp} % provide euro and other symbols
\else % if luatex or xetex
  \usepackage{unicode-math} % this also loads fontspec
  \defaultfontfeatures{Scale=MatchLowercase}
  \defaultfontfeatures[\rmfamily]{Ligatures=TeX,Scale=1}
\fi
\usepackage{lmodern}
\ifPDFTeX\else
  % xetex/luatex font selection
  \setmainfont[]{DejaVu Serif}
  \setmonofont[]{DejaVu Sans Mono}
\fi
% Use upquote if available, for straight quotes in verbatim environments
\IfFileExists{upquote.sty}{\usepackage{upquote}}{}
\IfFileExists{microtype.sty}{% use microtype if available
  \usepackage[]{microtype}
  \UseMicrotypeSet[protrusion]{basicmath} % disable protrusion for tt fonts
}{}
\usepackage{setspace}
\makeatletter
\@ifundefined{KOMAClassName}{% if non-KOMA class
  \IfFileExists{parskip.sty}{%
    \usepackage{parskip}
  }{% else
    \setlength{\parindent}{0pt}
    \setlength{\parskip}{6pt plus 2pt minus 1pt}}
}{% if KOMA class
  \KOMAoptions{parskip=half}}
\makeatother
\usepackage{listings}
\newcommand{\passthrough}[1]{#1}
\lstset{defaultdialect=[5.3]Lua}
\lstset{defaultdialect=[x86masm]Assembler}
\usepackage{graphicx}
\makeatletter
\newsavebox\pandoc@box
\newcommand*\pandocbounded[1]{% scales image to fit in text height/width
  \sbox\pandoc@box{#1}%
  \Gscale@div\@tempa{\textheight}{\dimexpr\ht\pandoc@box+\dp\pandoc@box\relax}%
  \Gscale@div\@tempb{\linewidth}{\wd\pandoc@box}%
  \ifdim\@tempb\p@<\@tempa\p@\let\@tempa\@tempb\fi% select the smaller of both
  \ifdim\@tempa\p@<\p@\scalebox{\@tempa}{\usebox\pandoc@box}%
  \else\usebox{\pandoc@box}%
  \fi%
}
% Set default figure placement to htbp
\def\fps@figure{htbp}
\makeatother
\setlength{\emergencystretch}{3em} % prevent overfull lines
\providecommand{\tightlist}{%
  \setlength{\itemsep}{0pt}\setlength{\parskip}{0pt}}
% Enable graphics inclusion and ensure figure numbering works
\usepackage{graphicx}
\renewcommand{\figurename}{Figure}

% Configure fonts for Unicode support with fallbacks
\usepackage{newunicodechar}
\newunicodechar{⁴}{\textsuperscript{4}}
\newunicodechar{₄}{\textsubscript{4}}

% Math notation macros
\newcommand{\norm}[1]{\lVert #1\rVert}
\newcommand{\abs}[1]{\lvert #1\rvert}

% Enhanced code block styling for better contrast and readability
\usepackage{fancyvrb}
\usepackage{xcolor}
\usepackage{listings}

% Define custom colors for code blocks
\definecolor{codebg}{RGB}{245, 245, 245}      % Light gray background
\definecolor{codeborder}{RGB}{200, 200, 200}  % Medium gray border
\definecolor{codefg}{RGB}{50, 50, 50}         % Dark gray text

% Configure Verbatim environment for inline code
\DefineVerbatimEnvironment{Verbatim}{Verbatim}{%
    fontsize=\small,
    frame=single,
    framerule=0.5pt,
    framesep=3pt,
    rulecolor=\color{codeborder},
    bgcolor=\color{codebg},
    fgcolor=\color{codefg}
}

% Configure code block styling
\DefineVerbatimEnvironment{Highlighting}{Verbatim}{%
    fontsize=\footnotesize,
    frame=single,
    framerule=0.5pt,
    framesep=5pt,
    rulecolor=\color{codeborder},
    bgcolor=\color{codebg},
    fgcolor=\color{codefg}
}

% Style inline code with \texttt
\renewcommand{\texttt}[1]{%
    \colorbox{codebg}{\color{codefg}\ttfamily #1}%
}

% Configure listings package for code blocks
\lstset{
    backgroundcolor=\color{codebg},
    basicstyle=\footnotesize\ttfamily\color{codefg},
    breakatwhitespace=false,
    breaklines=true,
    captionpos=b,
    commentstyle=\color{codefg},
    deletekeywords={...},
    escapeinside={\%*}{*)},
    extendedchars=true,
    frame=single,
    framerule=0.5pt,
    framesep=5pt,
    keepspaces=true,
    keywordstyle=\color{codefg},
    language=Python,
    morekeywords={*,...},
    numbers=left,
    numbersep=5pt,
    numberstyle=\tiny\color{codefg},
    rulecolor=\color{codeborder},
    showspaces=false,
    showstringspaces=false,
    showtabs=false,
    stepnumber=1,
    stringstyle=\color{codefg},
    tabsize=2,
    title=\lstname
}

% Override any Pandoc default lstset configurations
\AtBeginDocument{
    \lstset{
        backgroundcolor=\color{codebg},
        basicstyle=\footnotesize\ttfamily\color{codefg},
        frame=single,
        framerule=0.5pt,
        framesep=5pt,
        rulecolor=\color{codeborder},
        numbers=left,
        numbersep=5pt,
        numberstyle=\tiny\color{codefg}
    }
}

% Configure hyperref colors consistently
\AtBeginDocument{
% Override pandoc's hidelinks setting with consistent options
\hypersetup{
    colorlinks=true,
    allcolors=red,
    linkcolor=red,
    urlcolor=red,
    citecolor=red,
    filecolor=red,
    menucolor=red,
    linktoc=all
}
}

% Simple page break support for document structure
\usepackage{bookmark}
\IfFileExists{xurl.sty}{\usepackage{xurl}}{} % add URL line breaks if available
\urlstyle{same}
% fallback for those not using the hyperref driver hyperxmp:
\makeatletter
\@ifundefined{xmpquote}{\newcommand{\xmpquote}[1]{#1}}{}
\makeatother
\hypersetup{
  colorlinks=true,
  linkcolor={red},
  filecolor={red},
  citecolor={red},
  urlcolor={red},
  pdfcreator={LaTeX via pandoc}}

\title{Benchmarks and Statistics}
\author{Daniel Ari Friedman\\ ORCID: 0000-0001-6232-9096\\ Email: daniel@activeinference.institute\\ DOI: 10.5281/zenodo.16887791}
\date{October 07, 2026}

\begin{document}
\maketitle

{
\hypersetup{linkcolor=black}
\setcounter{tocdepth}{3}
\tableofcontents
}
\setstretch{1.0}
\section{Benchmarks and Statistics}\label{benchmarks-and-statistics}

\subsection{Overview}\label{overview}

This section documents the measurement surface behind the preceding
chapters. It is split across three modules with one job each:
\passthrough{\lstinline!src/quadmath/stats/benchmarks.py!} measures (a
small \passthrough{\lstinline!perf\_counter!} timing harness over the
core numerical surfaces),
\passthrough{\lstinline!src/quadmath/stats/statistics.py!} analyzes
(resampling inference, effect sizes, multiplicity correction, and
scaling fits, on standalone numpy), and
\passthrough{\lstinline!src/quadmath/viz/vis\_stats.py!} renders
(input-agnostic matplotlib primitives; the figures they produce are
collected in \passthrough{\lstinline!18\_stats\_gallery.md!}). The
measured workloads are the four surfaces the manuscript already treats:
quadray conversions (\passthrough{\lstinline!quadray.py!}, conventions
in \passthrough{\lstinline!SPEC.md!}), shell enumeration
(\passthrough{\lstinline!omni\_numbering.py!},
\passthrough{\lstinline!13\_lattice\_tooling.md!}), nearest-site search
(\passthrough{\lstinline!lattice\_search.py!}), and field fitting
(\passthrough{\lstinline!ivm\_field.py!},
\passthrough{\lstinline!11\_ivm\_field\_learning.md!}). Everything here
is deterministic: every random draw comes from a seeded
\passthrough{\lstinline!numpy.random.default\_rng!}, and every timing
routine is fully specified by its arguments, so the numbers below
reproduce bit-for-bit.

\subsection{Timing methodology}\label{timing-methodology}

\passthrough{\lstinline!time\_callable(fn, *, trials=5, warmup=1)!}
measures wall-clock duration with
\passthrough{\lstinline!time.perf\_counter!}, the highest-resolution
monotonic clock available. It runs one warmup call first and discards it
(first-call effects: module imports, key caches), then times
\passthrough{\lstinline!trials!} calls and reports the summary
statistics

\begin{equation}
\label{eq:bench-mean}
\bar{t} \;=\; \frac{1}{T}\sum_{i=1}^{T} t_i ,
\qquad t_m \;=\; \operatorname{median}(t_1,\dots,t_T),
\end{equation}

\begin{equation}
\label{eq:bench-percentile}
t_{95} \;=\; Q_{0.95}(t_1,\dots,t_T),
\end{equation}

where \(T\) is the number of trials and \(Q_p\) the empirical
\(p\)-quantile. The median is robust to a single outlier trial; the 95th
percentile bounds the tail that a user would actually feel. Results are
returned as \passthrough{\lstinline!BenchRow!}, a frozen dataclass with
fields \passthrough{\lstinline!name!}, \passthrough{\lstinline!n!},
\passthrough{\lstinline!trials!}, \passthrough{\lstinline!total\_s!},
\passthrough{\lstinline!mean\_s!}, \passthrough{\lstinline!median\_s!},
\passthrough{\lstinline!p95\_s!}, and
\passthrough{\lstinline!ops\_per\_s!}, plus
\passthrough{\lstinline!as\_dict()!} for plain data. Here
\passthrough{\lstinline!n!} is the per-call workload size (conversions
per call, shell depth, sites or queries per call), and throughput is

\begin{equation}
\label{eq:bench-ops}
\text{ops\_per\_s} \;=\; \frac{n}{\bar{t}} .
\end{equation}

Four constructors wrap the core surfaces with these defaults:
\passthrough{\lstinline!bench\_conversions(n=200, trials=5)!} times
round-trip quadray/embedding conversions,
\passthrough{\lstinline!bench\_shell\_enumeration(k\_max=4, trials=5)!}
times shell enumeration through
\passthrough{\lstinline!omni\_numbering.generate\_shell!},
\passthrough{\lstinline!bench\_lattice\_search(n\_sites=200, queries=50, trials=5)!}
times \passthrough{\lstinline!lattice\_search!} nearest-site queries
over a random ball of sites, and
\passthrough{\lstinline!bench\_field\_fit(n\_sites=64, trials=3)!} times
a synthetic field fit through \passthrough{\lstinline!IVMField.learn!}.
\passthrough{\lstinline!run\_all()!} executes all four with the shared
defaults collected in the module constant
\passthrough{\lstinline!BENCH\_DEFAULTS!} and returns the list of
\passthrough{\lstinline!BenchRow!} records;
\passthrough{\lstinline!summary\_table(rows)!} renders them as a
fixed-width text table. The harness measures only; it asserts nothing
--- timing distributions belong to analysis, not to test assertions.

\subsection{Percentile bootstrap}\label{percentile-bootstrap}

\passthrough{\lstinline!bootstrap\_ci(x, stat=np.mean, *, iters=2000, seed=0, alpha=0.05)!}
gives a confidence interval for any statistic \(\hat{\theta}\) by
resampling the data with replacement. With \(n = |x|\) and a seeded
\passthrough{\lstinline!numpy.random.default\_rng(seed)!}, iteration
\(b\) draws resampling indices
\(i_{bj} \sim \mathrm{Uniform}\{0,\dots,n-1\}\) in one vectorized call
(\passthrough{\lstinline!rng.integers(0, n, size=(iters, n))!}) and
evaluates

\begin{equation}
\label{eq:stat-bootstrap}
\hat{\theta}^{*}_{b} \;=\; \hat{\theta}\big(x_{i_{b1}},\dots,x_{i_{bn}}\big),
\qquad
\mathrm{CI}_{1-\alpha}
\;=\;
\Big[\, Q_{\alpha/2}\big(\hat{\theta}^{*}_{1..B}\big),\;
Q_{1-\alpha/2}\big(\hat{\theta}^{*}_{1..B}\big) \Big],
\end{equation}

with \(B\) = \passthrough{\lstinline!iters!}. The percentile interval
needs no distributional assumption about \(\hat{\theta}\) --- only that
the empirical resampling distribution approximates its sampling
distribution. All randomness is consumed from the seeded generator, so
the interval is exactly reproducible.

\subsection{Pooled permutation tests}\label{pooled-permutation-tests}

\passthrough{\lstinline!permutation\_test(a, b, *, iters=2000, seed=0, alternative="two-sided")!}
tests whether two samples differ in location without assuming a
distribution. Both samples are pooled; the observed statistic is
\(o = \big|\bar{a} - \bar{b}\big|\); each of the
\passthrough{\lstinline!iters!} iterations draws
\passthrough{\lstinline!rng.permutation!} of the pool, splits it at
\(\lvert a\rvert\) into a fake \(a\) and \(b\), and recomputes the
difference. The two-sided p-value uses the add-one convention

\begin{equation}
\label{eq:stat-permutation}
p \;=\;
\frac{1 + \#\big\{ r :\; |d_r| \,\ge\, o \big\}}{\text{iters} + 1},
\end{equation}

where \(d_r\) is the difference of permutation \(r\). The \(+1\) in
numerator and denominator counts the observed arrangement itself and
guarantees \(p \ge 1/(\text{iters}+1) > 0\): a permutation test can
never report an impossible zero p-value, and the estimate is
conservative. \passthrough{\lstinline!alternative="greater"!} and
\passthrough{\lstinline!"less"!} use the same convention with signed
comparisons of \(d_r\) against the signed observed difference.

\subsection{Effect size, multiplicity, and
scaling}\label{effect-size-multiplicity-and-scaling}

\passthrough{\lstinline!cohens\_d(a, b)!} reports the standardized mean
difference

\begin{equation}
\label{eq:stat-cohens}
d \;=\; \frac{\bar{a} - \bar{b}}{s_p},
\qquad
s_p \;=\;
\sqrt{\frac{(n_a - 1)\,s_a^2 + (n_b - 1)\,s_b^2}{n_a + n_b - 2}},
\end{equation}

with \(s_a, s_b\) the unbiased sample standard deviations. A p-value
says whether a difference is distinguishable from noise; \(d\) says how
large the difference is, which is the question that matters when
comparing benchmark configurations.

\passthrough{\lstinline!p\_adjust\_bonferroni(pvals)!} controls the
family-wise error rate when \(m\) hypotheses are tested at once:

\begin{equation}
\label{eq:stat-bonferroni}
p'_i \;=\; \min\!\big(1,\; m\, p_i\big).
\end{equation}

Each \(p'_i\) can be read as a valid p-value at level \(\alpha\) for the
whole family --- conservative, but assumption-free.

\passthrough{\lstinline!scaling\_fit(sizes, times)!} fits ordinary least
squares on the log-log data,

\begin{equation}
\label{eq:stat-scaling}
\log t \;=\; \beta_1 \log n + \beta_0,
\end{equation}

and returns \passthrough{\lstinline!(slope, intercept, r2)!}. The slope
is the empirical complexity exponent: \(\beta_1 \approx 1\) indicates
linear work in the input size \(n\), \(\beta_1 \approx 2\) quadratic,
and the coefficient of determination \(r^2\) measures how well the power
law describes the measurements.

\subsection{Jackknife, FDR, Welch, and circular
statistics}\label{sec:jackknife_fdr_welch_circular}

\passthrough{\lstinline!jackknife\_ci(x, stat=np.mean, alpha=0.05)!}
produces both a confidence interval and a bias estimate for any scalar
statistic without consuming a single random draw. Where the percentile
bootstrap of \eqref{eq:stat-bootstrap} resamples with replacement, the
jackknife deletes one observation at a time: each of the \(n\)
leave-one-out samples is scored, and the spread of those scores
determines the interval. With \(\hat{\theta}\) the full-sample statistic
and \(\hat{\theta}_{(i)}\) the statistic on the sample with observation
\(i\) deleted, the returned triple
\passthrough{\lstinline!(low, high, bias)!} is

\begin{equation}
\label{eq:stat-jackknife}
\begin{aligned}
\mathrm{bias} &= (n-1)\big(\bar{\theta}_{(\cdot)} - \hat{\theta}\big), &
\hat{\theta}_{\text{corr}} &= \hat{\theta} - \mathrm{bias}, \\
\mathrm{se}_{\text{jack}} &=
\sqrt{\tfrac{n-1}{n} \sum_{i=1}^{n} \big(\hat{\theta}_{(i)} -
\bar{\theta}_{(\cdot)}\big)^{2}}, &
\mathrm{CI}_{1-\alpha} &=
\hat{\theta} \pm z_{1-\alpha/2}\,\mathrm{se}_{\text{jack}},
\end{aligned}
\end{equation}

where \(\bar{\theta}_{(\cdot)}\) is the mean of the leave-one-out
scores. The bias-corrected estimate \(\hat{\theta}_{\text{corr}}\)
removes the first-order bias that the jackknife detects in curved
statistics; the interval itself stays centered on the uncorrected
\(\hat{\theta}\). The normal quantile \(z\) comes from bisecting
\passthrough{\lstinline!math.erfc(z / sqrt(2))!}, which decreases
monotonically from \(2\) to \(0\) as \(z\) runs from \(-\infty\) to
\(\infty\), so one hundred halvings of the bracket \([-40, 40]\) pin
\(z\) to double precision. This deterministic erfc-bisection quantile
needs no seed and no lookup table, so repeated calls with identical
arguments return bit-identical results --- the same reproducibility
contract as the seeded routines above, but with no generator to seed at
all.

\passthrough{\lstinline!benjamini\_hochberg(pvals)!} controls a
different error rate from the Bonferroni map of
\eqref{eq:stat-bonferroni}: instead of the family-wise error rate it
controls the expected proportion of false discoveries among rejections,
the false discovery rate. The step-up procedure sorts the \(m\) p-values
ascending, forms the raw values \(m\,p_{(i)} / i\) for ranks
\(i = 1,\dots,m\), and then enforces monotonicity with a running minimum
taken from the largest rank backwards,

\begin{equation}
\label{eq:stat-bh}
p'_{(i)} \;=\;
\min\!\Big(1,\; \min_{j \ge i}\, \frac{m\,p_{(j)}}{j}\Big),
\qquad
p'_{(m)} \;=\; \min(1,\, p_{(m)}),
\end{equation}

so the adjusted values are monotone non-decreasing in the sorted
p-values and never exceed the largest raw p-value. The backward
cumulative minimum is the step that the naive \(m\,p_{(i)}/i\) mapping
misses: without it, a large early-rank raw value could exceed a smaller
later-rank one, breaking monotonicity and invalidating the FDR
guarantee. The stable argsort of the input is recorded so the adjusted
values can be mapped back to the original input order, and the final cap
at \(1.0\) keeps the output a valid p-value vector.

\passthrough{\lstinline!welch\_t\_test(a, b, alternative="two-sided")!}
compares two sample means without assuming equal variances, the
assumption that the pooled standard deviation \(s_p\) of
\passthrough{\lstinline!cohens\_d!} makes. Writing \(w_a = s_a^2/n_a\)
and \(w_b = s_b^2/n_b\), the statistic and its Welch--Satterthwaite
degrees of freedom are

\begin{equation}
\label{eq:stat-welch}
t \;=\; \frac{\bar{a} - \bar{b}}{\sqrt{w_a + w_b}},
\qquad
\nu \;=\;
\frac{(w_a + w_b)^{2}}
{\dfrac{w_a^{2}}{n_a - 1} + \dfrac{w_b^{2}}{n_b - 1}} .
\end{equation}

The p-value is evaluated from the exact Student-\(t\) tails rather than
a normal approximation, and it needs no \passthrough{\lstinline!scipy!}:
the two-sided tail uses the identity

\begin{equation}
\label{eq:stat-tail}
\mathrm{P}\big(|T| \ge |t|\big) \;=\; I_{x}\!\big(\tfrac{\nu}{2},
\tfrac{1}{2}\big),
\qquad
x \;=\; \frac{\nu}{\nu + t^{2}},
\end{equation}

where \(I_x(a, b)\) is the regularized incomplete beta function computed
by a Lentz continued fraction (with the standard symmetry swap above
\(x = (a+1)/(a+b+2)\), an \(\exp\)-\(\log\Gamma\) front factor, and
guard floors against vanishing denominators).
\passthrough{\lstinline!"greater"!} and \passthrough{\lstinline!"less"!}
reuse the same two-sided tail, halved on the appropriate side of
\(t = 0\). The degenerate branches matter for constant inputs: when both
samples are constant the denominator of \eqref{eq:stat-welch} vanishes,
and equal means give \(t = 0.0\) (p-value \(1.0\) two-sided) while
different means give a signed infinite \(t\) whose tail probabilities
are exactly \(0\) or \(1\), evaluated against the finite pooled fallback
\(\nu = n_a + n_b - 2\) so the tails stay well-defined.

\passthrough{\lstinline!rotation\_stats(angles)!} treats angles in
radians as points on a circle rather than points on a line, where
\(359°\) and \(1°\) are neighbors:

\begin{equation}
\label{eq:stat-circular}
\bar{\theta} \;=\; \operatorname{atan2}\!\Big(
\tfrac{1}{n}\textstyle\sum_{j} \sin\theta_{j},\;
\tfrac{1}{n}\textstyle\sum_{j} \cos\theta_{j}\Big),
\qquad
R \;=\; \Big|\tfrac{1}{n}\textstyle\sum_{j} e^{i\theta_{j}}\Big| .
\end{equation}

The circular mean is the \passthrough{\lstinline!atan2!} of the averaged
sine and cosine, wrapped into \((-\pi, \pi]\), and the mean resultant
length \(R \in [0, 1]\) measures concentration: \(R = 1\) only when
every angle agrees, \(R = 0\) when the unit vectors cancel completely.
The returned dictionary gives the circular variance \(1 - R\) in
\([0, 1]\), and \passthrough{\lstinline!variance\_2pi!} \(= 2\,(1 - R)\)
under the convention for angles on the full \([0, 2\pi)\) circle, which
ranges over \([0, 2]\) and approaches the familiar linear variance for
tightly clustered angles. Because sine and cosine are \(2\pi\)-periodic,
wrapping the angles into any full circle leaves every returned value
unchanged; when the resultant vanishes the mean degenerates to whatever
\passthrough{\lstinline!atan2!} returns for the cancelled components.
Like everything else in
\passthrough{\lstinline!src/quadmath/stats/statistics.py!}, all four
routines are pure
\passthrough{\lstinline!numpy!}-plus-\passthrough{\lstinline!math!}
computations: no randomness, no optional scientific stack, and
bit-identical repeats for identical arguments.

\subsection{A reproducible example}\label{a-reproducible-example}

Both workhorses above are fully determined by their arguments, so the
following numbers are exact for any reader. The first call computes the
95\% bootstrap confidence interval of the mean of the first seven
Fibonacci numbers; the second asks whether the shifted ranges \(1..10\)
and \(11..20\) differ in location:

\begin{lstlisting}[language=Python]
import numpy as np

from quadmath.stats.statistics import bootstrap_ci, permutation_test

x = [2, 3, 5, 8, 13, 21, 34]
lo, hi = bootstrap_ci(x, np.mean, iters=2000, seed=0, alpha=0.05)
# (5.142857142857143, 20.571428571428573)

a = list(range(1, 11))
b = list(range(11, 21))
p = permutation_test(a, b, iters=2000, seed=0, alternative="two-sided")
# 0.0014992503748125937
\end{lstlisting}

With
\passthrough{\lstinline!bootstrap\_ci(x, np.mean, iters=2000, seed=0, alpha=0.05)!}
the mean of \(x\) is \(86/7 \approx 12.285714\) and the 95\% percentile
bootstrap interval is \([5.142857142857143,\; 20.571428571428573]\).
Both endpoints are exact multiples of \(1/7\): the mean of any 7-value
resample with repetition is a multiple of \(1/7\), so the percentile
grid is discrete rather than dense. With
\passthrough{\lstinline!permutation\_test(a, b, iters=2000, seed=0, alternative="two-sided")!}
on \(a = 1,\dots,10\) and \(b = 11,\dots,20\), the observed difference
is \(o = |\bar{a} - \bar{b}| = 10\) exactly, and exactly 2 of the 2000
seeded permutations reach it, giving

\begin{equation}
\label{eq:stat-example-p}
p \;=\; \frac{2 + 1}{2000 + 1} \;=\; \frac{3}{2001}
\;\approx\; 0.00149925 .
\end{equation}

The count of 2 deserves a comment. Among the
\(\binom{20}{10} = 184{,}756\) possible splits of the pool, exactly two
achieve \(|d_r| = 10\) --- the observed arrangement (\(1..10\) against
\(11..20\)) and its complement --- so the expected hit count over 2000
random draws is only \(2000 \cdot 2/184756 \approx 0.02\). The add-one
floor of \eqref{eq:stat-permutation} is \(1/2001 \approx 0.0005\), which
is what a typical seed reports here; the stream seeded with
\passthrough{\lstinline!seed=0!} happens to draw both exact-maximum
arrangements (at iterations 70 and 1146 of the 2000), so the count is 2
and the reported p-value is \(3/2001\). The example illustrates the
convention rather than a typical run: with any other seed the same call
reproduces the floor value, and in either case the test correctly never
reports the impossible \(p = 0\).

The two interval families meet on common ground in a second example. Two
synthetic lattice-error samples with a planted mean shift are drawn once
from a seeded generator: routine A contributes \(n = 64\) error
residuals centered on \(0.0\) and routine B \(64\) residuals whose
population mean sits \(0.08\) higher, both with scale \(0.15\) --- the
unit being whatever the error residual measures. For each sample mean,
the 95\% percentile bootstrap interval of \eqref{eq:stat-bootstrap}
(\(B = 2000\) resamples, re-seeded at 17) and the jackknife interval of
\eqref{eq:stat-jackknife} (which consumes no randomness at all) are
computed, and the script
\passthrough{\lstinline!quadmath/scripts/stats\_diagnostics\_gallery.py!}
draws all four intervals in a single
\passthrough{\lstinline!plot\_ci\_bars!} call:

\begin{lstlisting}[language=Python]
import numpy as np

from quadmath.stats.statistics import bootstrap_ci, jackknife_ci

rng = np.random.default_rng(17)
a = rng.normal(0.0, 0.15, size=64)
b = rng.normal(0.08, 0.15, size=64)
bootstrap_ci(a, np.mean, iters=2000, seed=17, alpha=0.05)
# (-0.07609532416913227, -0.0004865391750512773)
jackknife_ci(a, np.mean, alpha=0.05)[:2]
# (-0.07682589122849623, 0.0013292867175937334)
\end{lstlisting}

What the two families assume is where they differ. The percentile
bootstrap needs no shape assumption at all --- only that the empirical
resampling distribution of \eqref{eq:stat-bootstrap} approximates the
sampling distribution of \(\hat{\theta}\) --- while the jackknife
interval is the normal approximation \(\hat{\theta} \pm
z_{1-\alpha/2}\,\mathrm{se}_{\text{jack}}\) built from the \(n\)
leave-one-out scores of \eqref{eq:stat-jackknife}. On the sample mean
with \(n = 64\) the two agree to within a whisker: routine A's bootstrap
interval is \([-0.076095,\,-0.000487]\) against the jackknife's
\([-0.076826,\,0.001329]\), and routine B's is \([0.056788,\,0.132418]\)
against \([0.056011,\,0.130639]\) (sample means \(-0.037748\) and
\(0.093325\); the jackknife bias estimates sit at floating-point zero,
as \eqref{eq:stat-jackknife} predicts for a linear statistic). The
figure records the comparison: routine B's planted shift shows up as
both of its intervals clearing the dashed zero line, while routine A
hugs the line so tightly that the bootstrap upper endpoint stops
\(0.0005\) short of it and the jackknife upper endpoint crosses it by
\(0.0013\) --- a hair-width disagreement that is the honest lesson of
the panel. Coverage verdicts this close to the boundary are
method-sensitive; a Welch test on the same samples reports
\(t \approx -4.75\) with \(p \approx 5.4 \times 10^{-6}\)
(\eqref{eq:stat-welch}), in agreement with the shift both interval
families detect.

\begin{figure}
\centering
\pandocbounded{\includegraphics[keepaspectratio,alt={Bootstrap versus jackknife 95\% confidence intervals for two lattice-error samples. Sample means (dots) of two seeded synthetic error samples --- n = 64 per group from numpy.random.default\_rng(17), scale 0.15, planted mean shift 0.08 in the B population --- each carrying a 95\% percentile bootstrap interval (bootstrap\_ci, 2000 resamples, seed 17, \textbackslash alpha = 0.05) and a deterministic normal-approximation jackknife interval (jackknife\_ci, \textbackslash alpha = 0.05), drawn together in one plot\_ci\_bars call; the dashed line marks the unbiased population mean 0.0. Routine B's intervals exclude 0 while routine A's straddle it, and the two CI families agree within whisker width on both samples. Regenerate: uv run python quadmath/scripts/stats\_diagnostics\_gallery.py.}]{../output/figures/stats_ci_comparison.png}}
\caption{\textbf{Bootstrap versus jackknife 95\% confidence intervals
for two lattice-error samples.} Sample means (dots) of two seeded
synthetic error samples --- \(n = 64\) per group from
\protect\passthrough{\lstinline!numpy.random.default\_rng(17)!}, scale
0.15, planted mean shift 0.08 in the B population --- each carrying a
95\% percentile bootstrap interval
(\protect\passthrough{\lstinline!bootstrap\_ci!}, 2000 resamples, seed
17, \(\alpha = 0.05\)) and a deterministic normal-approximation
jackknife interval (\protect\passthrough{\lstinline!jackknife\_ci!},
\(\alpha = 0.05\)), drawn together in one
\protect\passthrough{\lstinline!plot\_ci\_bars!} call; the dashed line
marks the unbiased population mean 0.0. Routine B's intervals exclude 0
while routine A's straddle it, and the two CI families agree within
whisker width on both samples. Regenerate:
\protect\passthrough{\lstinline!uv run python quadmath/scripts/stats\_diagnostics\_gallery.py!}.}
\end{figure}

\subsection{Cross-references}\label{cross-references}

\begin{itemize}
\tightlist
\item
  Timing harness, \passthrough{\lstinline!BenchRow!}, and the four
  benchmark constructors:
  \passthrough{\lstinline!src/quadmath/stats/benchmarks.py!}.
\item
  Bootstrap, permutation tests, \passthrough{\lstinline!cohens\_d!},
  \passthrough{\lstinline!p\_adjust\_bonferroni!},
  \passthrough{\lstinline!scaling\_fit!}:
  \passthrough{\lstinline!src/quadmath/stats/statistics.py!}.
\item
  The figure primitives built on these results:
  \passthrough{\lstinline!18\_stats\_gallery.md!}.
\item
  The interval-comparison script and figure behind the second example:
  \passthrough{\lstinline!quadmath/scripts/stats\_diagnostics\_gallery.py!}
  (\passthrough{\lstinline!quadmath/output/figures/stats\_ci\_comparison.png!}).
\item
  Measured surfaces: \passthrough{\lstinline!quadray.py!} (conventions
  in \passthrough{\lstinline!SPEC.md!}),
  \passthrough{\lstinline!omni\_numbering.py!} and
  \passthrough{\lstinline!lattice\_search.py!}
  (\passthrough{\lstinline!13\_lattice\_tooling.md!}),
  \passthrough{\lstinline!ivm\_field.py!}
  (\passthrough{\lstinline!11\_ivm\_field\_learning.md!}).
\item
  Sibling gallery of the lattice layer:
  \passthrough{\lstinline!16\_lattice\_gallery.md!}.
\end{itemize}

\end{document}
